% Options for packages loaded elsewhere
\PassOptionsToPackage{unicode}{hyperref}
\PassOptionsToPackage{hyphens}{url}
\documentclass[
]{article}
\usepackage{xcolor}
\usepackage{amsmath,amssymb}
\setcounter{secnumdepth}{5}
\usepackage{iftex}
\ifPDFTeX
  \usepackage[T1]{fontenc}
  \usepackage[utf8]{inputenc}
  \usepackage{textcomp} % provide euro and other symbols
\else % if luatex or xetex
  \usepackage{unicode-math} % this also loads fontspec
  \defaultfontfeatures{Scale=MatchLowercase}
  \defaultfontfeatures[\rmfamily]{Ligatures=TeX,Scale=1}
\fi
% \usepackage{lmodern}
\ifPDFTeX\else
  % xetex/luatex font selection
\fi
% Use upquote if available, for straight quotes in verbatim environments
\IfFileExists{upquote.sty}{\usepackage{upquote}}{}
\IfFileExists{microtype.sty}{% use microtype if available
  \usepackage[]{microtype}
  \UseMicrotypeSet[protrusion]{basicmath} % disable protrusion for tt fonts
}{}
\makeatletter
\@ifundefined{KOMAClassName}{% if non-KOMA class
  \IfFileExists{parskip.sty}{%
    \usepackage{parskip}
  }{% else
    \setlength{\parindent}{0pt}
    \setlength{\parskip}{6pt plus 2pt minus 1pt}}
}{% if KOMA class
  \KOMAoptions{parskip=half}}
\makeatother
\usepackage{longtable,booktabs,array}
\newcounter{none} % for unnumbered tables
\usepackage{calc} % for calculating minipage widths
% Correct order of tables after \paragraph or \subparagraph
\usepackage{etoolbox}
\makeatletter
\patchcmd\longtable{\par}{\if@noskipsec\mbox{}\fi\par}{}{}
\makeatother
% Allow footnotes in longtable head/foot
\IfFileExists{footnotehyper.sty}{\usepackage{footnotehyper}}{\usepackage{footnote}}
\makesavenoteenv{longtable}
\setlength{\emergencystretch}{3em} % prevent overfull lines
\providecommand{\tightlist}{%
  \setlength{\itemsep}{0pt}\setlength{\parskip}{0pt}}
\usepackage[]{natbib}
\bibliographystyle{plainnat}
\makeatletter
\@ifpackageloaded{subcaption}{}{\usepackage{subcaption}}
\@ifpackageloaded{caption}{}{\usepackage{caption}}
\captionsetup[subfigure]{margin=0.5em}
\AtBeginDocument{%
\renewcommand*\figurename{Figure}
\renewcommand*\tablename{Table}
}
\AtBeginDocument{%
\renewcommand*\listfigurename{List of Figures}
\renewcommand*\listtablename{List of Tables}
}
\newcounter{pandoccrossref@subfigures@footnote@counter}
\newenvironment{pandoccrossrefsubfigures}{%
\setcounter{pandoccrossref@subfigures@footnote@counter}{0}
\begin{figure}\centering%
\gdef\global@pandoccrossref@subfigures@footnotes{}%
\DeclareRobustCommand{\footnote}[1]{\footnotemark%
\stepcounter{pandoccrossref@subfigures@footnote@counter}%
\ifx\global@pandoccrossref@subfigures@footnotes\empty%
\gdef\global@pandoccrossref@subfigures@footnotes{{##1}}%
\else%
\g@addto@macro\global@pandoccrossref@subfigures@footnotes{, {##1}}%
\fi}}%
{\end{figure}%
\addtocounter{footnote}{-\value{pandoccrossref@subfigures@footnote@counter}}
\@for\f:=\global@pandoccrossref@subfigures@footnotes\do{\stepcounter{footnote}\footnotetext{\f}}%
\gdef\global@pandoccrossref@subfigures@footnotes{}}
\@ifpackageloaded{float}{}{\usepackage{float}}
\floatstyle{ruled}
\@ifundefined{c@chapter}{\newfloat{codelisting}{h}{lop}}{\newfloat{codelisting}{h}{lop}[chapter]}
\floatname{codelisting}{Listing}
\newcommand*\listoflistings{\listof{codelisting}{List of Listings}}
\makeatother
\usepackage{bookmark}
\IfFileExists{xurl.sty}{\usepackage{xurl}}{} % add URL line breaks if available
\urlstyle{same}
% fallback for those not using the hyperref driver hyperxmp:
\makeatletter
\@ifundefined{xmpquote}{\newcommand{\xmpquote}[1]{#1}}{}
\makeatother
\hypersetup{
  pdftitle={Redacted Report Template: Disclosure Control and Release Audit},
  pdfauthor={Daniel Ari Friedman},
  colorlinks=true,linkcolor=red,urlcolor=red,citecolor=red,anchorcolor=red,filecolor=red,
  pdfcreator={LaTeX via pandoc}}

\title{Redacted Report Template: Disclosure Control and Release Audit}
\author{Daniel Ari Friedman}
\date{}

\IfFileExists{seqsplit.sty}{\usepackage{seqsplit}}{\newcommand{\seqsplit}[1]{#1}}
\protected\def\breakseq#1{\seqsplit{#1}}
\protected\def\breaktt#1{\begingroup\ttfamily\seqsplit{#1}\endgroup}
\title{Redacted Report Template: Disclosure Control and Release Audit}
\author{Daniel Ari Friedman\\\footnotesize{Active Inference Institute}\\\footnotesize{\texttt{daniel@activeinference.institute}}\\\footnotesize{\href{https://orcid.org/0000-0001-6232-9096}{ORCID: 0000-0001-6232-9096}} \\ \footnotesize{DOI: \href{https://doi.org/10.5281/zenodo.21298890}{10.5281/zenodo.21298890}}}
\date{2026-07-10}

\usepackage[left=0.16in,right=0.16in,top=0.48in,bottom=0.52in]{geometry}
\usepackage{xcolor}
\usepackage{graphicx}
\definecolor{redactedReportPage}{HTML}{FFFFFF}
\pagecolor{redactedReportPage}
\color{black}
\setlength{\parindent}{0pt}
\setlength{\parskip}{0.45em}

\begin{document}
\begin{titlepage}
\centering
\vspace*{2cm}
{\LARGE\bfseries Redacted Report Template: Disclosure Control and Release Audit\par}
\vspace{0.75em}
{\large A public exemplar for sanitized release workflows\par}
\vfill
\makeatletter
{\@author\par}
\vspace{1em}
{\@date\par}
\makeatother
\vfill
\end{titlepage}
\thispagestyle{empty}
\tableofcontents
\newpage

\section{Abstract}\label{abstract}

This exemplar demonstrates a complete disclosure-control pipeline for
sanitized public release reports. The methodology combines
classification-ceiling enforcement, source-protection validation,
mosaic-risk scoring, and TPM-backed sealed sidecars across a
sixteen-variant visual proof matrix. Four redaction styles---blackout,
whiteout, grayout, and blur---are rendered across four PDF
backgrounds---white, gray, black, and blur---yielding sixteen base proof
PDFs. Each receives nine steganographic security methods including
SHA-256/SHA-512 hash manifests, diagonal watermark overlays, footer
provenance stamps, invisible text, QR and Code128 barcodes, PDF Info and
XMP metadata, and embedded manifest attachments. Optional Kmyth TPM
sealing wraps each hash manifest and steganography PDF in a
\texttt{.ski} sidecar sealed against the TPM2-TSS storage hierarchy,
bound to PCR selections and policy or-values. The release gate requires
three reviewer roles---originator, classification reviewer, and release
authority---each providing a non-empty rationale. A source-safe
redaction ledger records SHA-256 hashes of each redacted span without
exposing source text, and a segment hash manifest compares source and
public SHA-256 digests for reproducible audit. The comprehensive release
packet combines sanitized text, audit findings, ledger, hashes, review
gate status, and paragraph-level audit tables into a single JSON-ready
export. This exemplar confirms that visual presentation choices remain
orthogonal to the release gate: the same source-safe decisions drive
every output variant.

\newpage

\section{Introduction}\label{introduction}

Disclosure control---the process of sanitizing classified or sensitive
information before public release---requires multiple layers of
validation to ensure that no source identities, operational details,
time-place selectors, or controlled dissemination markers leak into the
public record. This exemplar implements a reproducible pipeline that
combines text-level disclosure control with visual redaction proofing,
steganographic provenance embedding, and hardware-backed TPM sealing.

The pipeline operates on invented fixture data: synthetic segments with
classification levels ranging from UNCLASSIFIED through
TOP\_SECRET//SCI, synthetic redaction decisions with bounded reasons,
and synthetic reviewer records. No real source material is used. The
exemplar demonstrates the full release-review workflow:
classification-ceiling enforcement, redaction-span validation,
source-control coverage, mosaic-risk scoring, residual-marker detection,
review-gate evaluation, source-safe ledger generation, and TPM-sealed
sidecar production.

\subsection{Scope and Safety
Boundaries}\label{scope-and-safety-boundaries}

This exemplar is limited to lawful redaction, declassification support,
public-records release review, taxonomy normalization, source-safe
ledgers, reviewer approval gates, sanitized packet export, and
source-protection auditing. It does not provide targeting, collection,
evasion, or surveillance operational guidance. All committed examples
are invented fixtures.

\subsection{Contributions}\label{contributions}

\begin{enumerate}
\def\labelenumi{\arabic{enumi}.}
\tightlist
\item
  A classification taxonomy with SCI alias support and configurable
  public ceilings.
\item
  A release-audit engine that validates redaction spans, checks orphan
  decisions, enforces source-control coverage, and scores mosaic risk.
\item
  A source-safe redaction ledger that records SHA-256 hashes of redacted
  spans without exposing source text.
\item
  A visual proof matrix that enumerates four redaction styles across
  four PDF backgrounds, producing sixteen variant PDFs with identical
  source-safe decisions.
\item
  A steganography layer that post-processes each base PDF with nine
  security methods, producing provenance-enhanced companion PDFs with
  hash manifests.
\item
  Optional Kmyth TPM sealing that wraps each hash manifest and
  steganography PDF in \texttt{.ski} sidecars sealed against the
  TPM2-TSS storage hierarchy.
\item
  A comprehensive release packet that combines sanitized text, audit
  findings, ledger, hashes, review gate status, and paragraph audit
  tables.
\end{enumerate}

\section{Architecture}\label{architecture}

The pipeline architecture consists of two layers. Layer one performs
text-level disclosure control: each segment is audited against a public
classification ceiling, redaction spans are validated for non-overlap,
orphan decisions are flagged, and source-control coverage is enforced.
Layer two applies visual redaction treatments, steganographic provenance
overlays, and optional Kmyth TPM sidecar sealing.

\subsection{Layer One: Disclosure
Control}\label{layer-one-disclosure-control}

The disclosure-control engine operates on \breaktt{RedactionSegment}
objects, each carrying an identifier, classification level, text
content, and optional source controls. Redaction decisions are
\breaktt{RedactionDecision} records with bounded reasons:
\texttt{source\_identity}, \breaktt{operational\_detail},
\breaktt{time\_place\_selector}, \texttt{legal\_privilege}, and
\texttt{privacy}. The audit engine validates that:

\begin{itemize}
\tightlist
\item
  Redaction spans are non-overlapping and within text bounds.
\item
  Each segment above the public ceiling has at least one redaction
  decision.
\item
  Each source control has a corresponding \texttt{source\_identity}
  redaction.
\item
  No orphan decisions reference missing segments.
\item
  Residual markers (email, IP, coordinates, NOFORN, HUMINT, SIGINT) are
  absent from sanitized text.
\item
  The mosaic risk score---residual markers normalized by segment and
  pattern count---does not exceed the policy threshold.
\end{itemize}

\subsection{Layer Two: Visual Proof and
Steganography}\label{layer-two-visual-proof-and-steganography}

Visual proofing is parameterized separately from the release-audit text
path. Four redaction styles---blackout (solid black fill), whiteout
(white fill with gray text), grayout (mid-gray fill), and blur
(offset-rendered token)---are rendered across four PDF
backgrounds---white, gray, black, and blur (subdued text with blur
effect). The 4\texorpdfstring{\ensuremath{\times}}{x}4 matrix yields sixteen variant PDFs, each receiving the
same source-safe redaction decisions.

The steganography layer post-processes each base PDF with nine security
methods:

{\def\LTcaptype{none} % do not increment counter
\begin{longtable}[]{@{}
  >{\raggedright\arraybackslash}p{(\linewidth - 2\tabcolsep) * \real{0.4706}}
  >{\raggedright\arraybackslash}p{(\linewidth - 2\tabcolsep) * \real{0.5294}}@{}}
\toprule\noalign{}
\begin{minipage}[b]{\linewidth}\raggedright
Method
\end{minipage} & \begin{minipage}[b]{\linewidth}\raggedright
Purpose
\end{minipage} \\
\midrule\noalign{}
\endhead
\bottomrule\noalign{}
\endlastfoot
SHA-256/SHA-512 hash manifest & Cryptographic integrity verification \\
Diagonal watermark overlay & Visible provenance stamp \\
Footer provenance overlay & Page-level release metadata \\
First-page invisible text & Hidden identification marker \\
QR payload barcode & Machine-readable provenance \\
Code128 page barcode & Per-page tracking barcode \\
PDF Info metadata & Document-level metadata injection \\
XMP metadata & Standards-compliant metadata embedding \\
Embedded stego manifest attachment & Self-contained provenance
archive \\
\end{longtable}
}

\subsection{Layer Three: Kmyth TPM
Sealing}\label{layer-three-kmyth-tpm-sealing}

Kmyth TPM sealing wraps each hash manifest and steganography PDF in a
\texttt{.ski} sidecar sealed against the TPM2-TSS storage hierarchy. The
sealed objects are bound to PCR selections and policy or-values,
ensuring that unsealing requires the same platform configuration.

On macOS, which lacks a hardware TPM, a software TPM emulator (swtpm)
and an mssim-to-swtpm protocol proxy bridge the TPM2-TSS mssim TCTI to
swtpm's native socket protocol. The proxy:

\begin{itemize}
\tightlist
\item
  \textbf{Control channel}: intercepts MS simulator platform commands
  (POWER\_ON=1, NV\_ON=11, TPM\_SESSION\_END=20) and returns success,
  since swtpm's \breaktt{-\/-flags\ startup-clear} handles TPM
  initialization.
\item
  \textbf{Data channel}: strips the 9-byte mssim wire header (4B
  cmd\_type + 1B locality + 4B tpm\_size), forwards raw TPM commands to
  swtpm, and wraps responses back in mssim format.
\end{itemize}

The kmyth-seal binary was patched to call
\breaktt{Tss2\_Sys\_FlushContext} for the storage key handle before
freeing TPM2 resources. Without this patch, consecutive kmyth-seal
invocations exhaust the swtpm transient object slots (typically three),
causing ``out of memory for object contexts'' errors on the second seal.

\newpage

\section{Results}\label{results}

\subsection{Fixture Release Packet}\label{fixture-release-packet}

The fixture release packet contains fourteen segments spanning four
classification levels: UNCLASSIFIED (ten segments), CUI (one segment),
SECRET (two segments), and TOP\_SECRET (one segment). Four segments
carry source controls (drawn from the HUMINT, SIGINT, and IMINT
disciplines). Twenty-one redaction decisions are applied across those
four segments, using all five bounded reasons:
\texttt{source\_identity}, \breaktt{operational\_detail},
\breaktt{time\_place\_selector}, \texttt{legal\_privilege}, and
\texttt{privacy}.

The audit produces:

\begin{itemize}
\tightlist
\item
  \textbf{Releasable}: true (no error-level findings after redaction).
\item
  \textbf{Release safety score}: weighted by error count, warning count,
  and mosaic risk.
\item
  \textbf{Redaction coverage}: 1.0 (all sensitive segments have at least
  one decision).
\item
  \textbf{Mosaic risk score}: residual markers normalized by segment and
  pattern count.
\item
  \textbf{Findings}: warning-level findings for residual markers in
  sanitized text.
\end{itemize}

\subsection{Source-Safe Redaction
Ledger}\label{source-safe-redaction-ledger}

The redaction ledger records each decision with:

\begin{itemize}
\tightlist
\item
  A decision ID derived from the SHA-256 of the segment ID, span,
  reason, and replacement.
\item
  The segment ID, start, end, span length, reason, and replacement.
\item
  A \texttt{valid\_span} flag indicating whether the span falls within
  the segment text.
\item
  A \breaktt{source\_span\_sha256} hash of the redacted text---present
  only for valid spans, absent for invalid or orphan decisions.
\end{itemize}

The ledger never exposes source text. It provides reproducible audit
evidence without disclosure risk.

\subsection{Segment Hash Manifest}\label{segment-hash-manifest}

The hash manifest records, for each segment:

\begin{itemize}
\tightlist
\item
  \texttt{source\_sha256}: SHA-256 of the original segment text.
\item
  \texttt{public\_sha256}: SHA-256 of the redacted segment text.
\end{itemize}

This allows downstream verification that the public release matches the
audited version without comparing source text directly.

\subsection{Review Gate}\label{review-gate}

The release gate requires three reviewer roles: originator,
classification reviewer, and release authority. Each reviewer provides a
non-empty rationale. The fixture reviews all approve, yielding:

\begin{itemize}
\tightlist
\item
  \textbf{Approved}: true.
\item
  \textbf{Approval count}: 3.
\item
  \textbf{Required roles present}: classification\_reviewer, originator,
  release\_authority.
\item
  \textbf{Findings}: none.
\end{itemize}

The \breaktt{final\_release\_recommended} flag is true only when the
packet is releasable, the review gate is approved, and no blocking
warnings remain.

\subsection{Visual Proof Matrix}\label{visual-proof-matrix}

The development proof matrix produces sixteen base PDFs (four redaction
styles \texorpdfstring{\ensuremath{\times}}{x} four backgrounds) and, when steganography is enabled, sixteen
companion steganography PDFs with hash manifests. Each variant records:

\begin{itemize}
\tightlist
\item
  Base PDF filename, byte size, and SHA-256.
\item
  Steganography PDF filename, byte size, and SHA-256.
\item
  Hash manifest filename and SHA-256.
\item
  Kmyth sidecar count and filenames (when Kmyth is available).
\end{itemize}

\subsection{Kmyth TPM Sidecar
Production}\label{kmyth-tpm-sidecar-production}

When Kmyth tools are runnable and a TPM backend is available, each
variant produces two \texttt{.ski} sidecars:

\begin{enumerate}
\def\labelenumi{\arabic{enumi}.}
\tightlist
\item
  \texttt{\{variant\_id\}.hashes.json.ski} --- the hash manifest sealed
  against the TPM storage hierarchy.
\item
  \texttt{\{variant\_id\}\_steganography.pdf.ski} --- the steganography
  PDF sealed against the TPM storage hierarchy.
\end{enumerate}

The \texttt{.ski} files are ASCII-armored and contain PCR selections,
policy or-values, the storage key public area, and the sealed data
object. Unsealing requires the same TPM and platform configuration.

In the verified run, all sixteen variants produced both sidecars,
yielding thirty-two \texttt{.ski} files total. The kmyth-seal binary's
FlushContext patch ensured that consecutive seal invocations did not
exhaust the swtpm transient object slots.

\subsection{Residual Risk Detection}\label{residual-risk-detection}

The residual-risk detector scans sanitized text for common
public-release leaks:

{\def\LTcaptype{none} % do not increment counter
\begin{longtable}[]{@{}
  >{\raggedright\arraybackslash}p{(\linewidth - 2\tabcolsep) * \real{0.5625}}
  >{\raggedright\arraybackslash}p{(\linewidth - 2\tabcolsep) * \real{0.4375}}@{}}
\toprule\noalign{}
\begin{minipage}[b]{\linewidth}\raggedright
Pattern
\end{minipage} & \begin{minipage}[b]{\linewidth}\raggedright
Regex
\end{minipage} \\
\midrule\noalign{}
\endhead
\bottomrule\noalign{}
\endlastfoot
Email address &
\texttt{\textbackslash{}b{[}A-Za-z0-9.\_\%+-{]}+@{[}A-Za-z0-9.-{]}+\textbackslash{}.{[}A-Za-z{]}\{2,\}\textbackslash{}b} \\
IPv4 address &
\texttt{\textbackslash{}b(?:\textbackslash{}d\{1,3\}\textbackslash{}.)\{3\}\textbackslash{}d\{1,3\}\textbackslash{}b} \\
Coordinate pair &
\texttt{\textbackslash{}b-?\textbackslash{}d\{1,2\}\textbackslash{}.\textbackslash{}d\{3,\},\textbackslash{}s*-?\textbackslash{}d\{1,3\}\textbackslash{}.\textbackslash{}d\{3,\}\textbackslash{}b} \\
Controlled dissemination &
\texttt{\textbackslash{}b(?:NOFORN\textbackslash{}|{}ORCON\textbackslash{}|{}REL\textbackslash{}s+TO)\textbackslash{}b} \\
Collection discipline &
\texttt{\textbackslash{}b(?:HUMINT\textbackslash{}|{}SIGINT\textbackslash{}|{}IMINT\textbackslash{}|{}MASINT\textbackslash{}|{}OSINT)\textbackslash{}b} \\
Compartment marker &
\texttt{\textbackslash{}b(?:SCI\textbackslash{}|{}TS\_SCI\textbackslash{}|{}TOP\textbackslash{}s+SECRET)\textbackslash{}b} \\
Sensitive markers & \texttt{HUMINT}, \texttt{SIGINT}, \texttt{source},
\texttt{selector}, \texttt{location}, \texttt{2026-} \\
\end{longtable}
}

Each detected marker generates a warning finding. The mosaic risk score
aggregates residual markers across all segments.

\section{Discussion}\label{discussion}

The exemplar demonstrates that disclosure control can be decomposed into
orthogonal concerns: text-level audit, visual presentation,
steganographic provenance, and hardware-backed sealing. Each concern is
independently configurable and verifiable.

The visual proof matrix confirms that blackout, whiteout, grayout, and
blur treatments produce equivalent source-safe outputs---only the visual
token differs. The steganography layer adds provenance without altering
the redaction decisions. Kmyth TPM sealing adds a hardware binding that
ensures sealed sidecars can only be unsealed on the same platform
configuration.

The mssim-to-swtpm protocol proxy is a necessary bridge on macOS, where
no hardware TPM exists. The proxy handles three protocol mismatches: (1)
platform commands use different command numbers, (2) the data channel
wraps TPM commands in a 9-byte mssim header, and (3) swtpm's transient
object slots are limited and must be flushed between seal invocations.

The FlushContext patch to kmyth-seal is a critical fix for batch sealing
workflows. Without it, the second kmyth-seal invocation fails with ``out
of memory for object contexts'' because the storage key from the first
invocation remains loaded in the TPM's transient object table.

\subsection{Limitations}\label{limitations}

\begin{itemize}
\tightlist
\item
  The fixture data is invented; real release packets may require
  additional classification taxonomies and review-role policies.
\item
  The swtpm software TPM does not provide hardware-level tamper
  resistance; production deployments should use a hardware TPM.
\item
  The mssim proxy adds latency to each TPM command; batch sealing of
  thirty-two sidecars takes approximately thirty seconds.
\item
  PDF password encryption is optional and uses AES-256; the password is
  not stored in the variant matrix.
\end{itemize}



\clearpage

\bibliography{references}
\end{document}
